\textbf{Phần 1. Trắc nghiệm lựa chọn}
\setcounterpageref{ex}{0}

\begin{ex} %[5H8N48]
Đơn vị đo thể tích nào sau đây lớn nhất?
\choice
{\True Mét khối}
{Xăng-ti-mét khối}
{Đề-xi-mét khối}
{Lít}
\loigiai{
Mét khối ($1\text{ m}^3 = 1\,000\text{ dm}^3 = 1\,000\,000\text{ cm}^3$) là đơn vị lớn nhất.\\
\textbf{Đáp án của câu hỏi là A}
}
\end{ex}

\begin{ex} %[5H8N48]
$3\text{ m}^3\; 45\text{ dm}^3$ được viết dưới dạng số thập phân có đơn vị mét khối là
\choice
{$3{,}45\text{ m}^3$}
{$3{,}450\text{ m}^3$}
{\True $3{,}045\text{ m}^3$}
{$34{,}5\text{ m}^3$}
\loigiai{
$3\text{ m}^3\; 45\text{ dm}^3 = 3 + \dfrac{45}{1\,000} = 3{,}045\text{ m}^3$.\\
\textbf{Đáp án của câu hỏi là C}
}
\end{ex}

\begin{ex} %[5H8H48]
Sắp xếp các số đo sau theo thứ tự từ bé đến lớn: $1{,}5\text{ m}^3$; $1\,200\text{ dm}^3$; $1\,600\,000\text{ cm}^3$.
\choice
{$1{,}5\text{ m}^3 < 1\,200\text{ dm}^3 < 1\,600\,000\text{ cm}^3$}
{$1\,600\,000\text{ cm}^3 < 1{,}5\text{ m}^3 < 1\,200\text{ dm}^3$}
{$1\,200\text{ dm}^3 < 1\,600\,000\text{ cm}^3 < 1{,}5\text{ m}^3$}
{\True $1\,200\text{ dm}^3 < 1{,}5\text{ m}^3 < 1\,600\,000\text{ cm}^3$}
\loigiai{
Đổi về $\text{dm}^3$: $1{,}5\text{ m}^3 = 1\,500\text{ dm}^3$; $1\,600\,000\text{ cm}^3 = 1\,600\text{ dm}^3$.\\
Thứ tự: $1\,200\text{ dm}^3 < 1{,}5\text{ m}^3 < 1\,600\,000\text{ cm}^3$.\\
\textbf{Đáp án của câu hỏi là D}
}
\end{ex}

\begin{ex} %[5H8H48]
Giá trị của biểu thức $0{,}75\text{ m}^3 + 350\text{ dm}^3 - 120\,000\text{ cm}^3$ theo đơn vị đề-xi-mét khối là
\choice
{$1\,100\text{ dm}^3$}
{\True $980\text{ dm}^3$}
{$880\text{ dm}^3$}
{$950\text{ dm}^3$}
\loigiai{
$750 + 350 - 120 = 980\text{ (dm}^3\text{)}$.\\
\textbf{Đáp án của câu hỏi là B}
}
\end{ex}

\begin{ex} %[5H8H48]
Một khối gỗ hình lập phương có thể tích $0{,}064\text{ m}^3$. Thể tích khối gỗ đó theo xăng-ti-mét khối là
\choice
{\True $64\,000\text{ cm}^3$}
{$640\text{ cm}^3$}
{$6\,400\text{ cm}^3$}
{$640\,000\text{ cm}^3$}
\loigiai{
$0{,}064 \times 1\,000\,000 = 64\,000\text{ cm}^3$.\\
\textbf{Đáp án của câu hỏi là A}
}
\end{ex}

\begin{ex} %[5H8H48]
Một bể nước chứa $2{,}4\text{ m}^3$ nước. Người ta đã dùng hết $40\%$ lượng nước trong bể. Lượng nước còn lại trong bể là
\choice
{$960\text{ l}$}
{$1\,200\text{ l}$}
{\True $1\,440\text{ l}$}
{$1\,600\text{ l}$}
\loigiai{
Lượng nước ban đầu: $2\,4\text{ m}^3 = 2\,400\text{ l}$.\\
Nước còn lại: $2\,400 \times 60\% = 1\,440\text{ l}$.\\
\textbf{Đáp án của câu hỏi là C}
}
\end{ex}

\begin{ex} %[5H8V48]
Một hình hộp chữ nhật có chiều dài $2\text{ m}$, chiều rộng $1{,}5\text{ m}$ và chiều cao $1{,}2\text{ m}$. Thể tích của hình hộp chữ nhật đó là
\choice
{$3\text{ m}^3$}
{$4{,}2\text{ m}^3$}
{$4{,}8\text{ m}^3$}
{\True $3{,}6\text{ m}^3$}
\loigiai{
$V = 2 \times 1{,}5 \times 1{,}2 = 3{,}6\text{ (m}^3\text{)}$.\\
\textbf{Đáp án của câu hỏi là D}
}
\end{ex}

\begin{ex} %[5H8V48]
Một bể nước hình hộp chữ nhật có diện tích đáy là $4{,}5\text{ m}^2$. Đang chứa $5\,400\text{ l}$ nước. Mực nước trong bể cao bao nhiêu mét?
\choice
{$1\text{ m}$}
{\True $1{,}2\text{ m}$}
{$1{,}5\text{ m}$}
{$0{,}8\text{ m}$}
\loigiai{
$5\,400\text{ l} = 5{,}4\text{ m}^3$.\\
Chiều cao mực nước: $5{,}4 : 4{,}5 = 1{,}2\text{ (m)}$.\\
\textbf{Đáp án của câu hỏi là B}
}
\end{ex}

\begin{ex} %[5H8V48]
Một khối kim loại hình lập phương cạnh $0{,}5\text{ m}$. Biết mỗi đề-xi-mét khối kim loại nặng $7{,}8\text{ kg}$. Khối lượng của khối kim loại đó là
\choice
{$97{,}5\text{ kg}$}
{$9\,750\text{ kg}$}
{\True $975\text{ kg}$}
{$390\text{ kg}$}
\loigiai{
Thể tích: $0{,}5 \times 0{,}5 \times 0{,}5 = 0{,}125\text{ m}^3 = 125\text{ dm}^3$.\\
Khối lượng: $125 \times 7{,}8 = 975\text{ (kg)}$.\\
\textbf{Đáp án của câu hỏi là C}
}
\end{ex}

\begin{ex} %[5H8V48]
Một bình tràn chứa đầy nước. Người ta thả vào bình một hòn đá thì có $2{,}5\text{ l}$ nước tràn ra. Thể tích của hòn đá đó là
\choice
{$25\text{ cm}^3$}
{\True $2\,500\text{ cm}^3$}
{$250\text{ cm}^3$}
{$0{,}25\text{ m}^3$}
\loigiai{
$2{,}5\text{ l} = 2{,}5\text{ dm}^3 = 2\,500\text{ cm}^3$.\\
\textbf{Đáp án của câu hỏi là B}
}
\end{ex}

\begin{ex} %[5H8V48]
Có hai bể nước hình hộp chữ nhật. Bể thứ nhất chứa $4\text{ m}^3\; 50\text{ dm}^3$ nước, bể thứ hai chứa $3\text{ m}^3\; 800\text{ dm}^3$ nước. Cả hai bể chứa tất cả số lít nước là
\choice
{\True $7\,850\text{ l}$}
{$7\,850\text{ l}$}
{$8\,300\text{ l}$}
{$7\,580\text{ l}$}
\loigiai{
Bể 1: $4\,050\text{ l}$; Bể 2: $3\,800\text{ l}$.\\
Tổng cộng: $4\,050 + 3\,800 = 7\,850\text{ (l)}$.\\
\textbf{Đáp án của câu hỏi là A}
}
\end{ex}

\begin{ex} %[5H8C48]
Một bể nước hình hộp chữ nhật có chiều dài $2{,}5\text{ m}$, chiều rộng $1{,}8\text{ m}$ và chiều cao $1{,}5\text{ m}$. Bể đang cạn nước, người ta mở vòi nước chảy vào bể mỗi phút được $45\text{ l}$. Sau bao lâu thì nước đầy $80\%$ bể?
\choice
{$100\text{ phút}$}
{$150\text{ phút}$}
{$90\text{ phút}$}
{\True $120\text{ phút}$}
\loigiai{
Thể tích bể: $2{,}5 \times 1{,}8 \times 1{,}5 = 6{,}75\text{ m}^3 = 6\,750\text{ l}$.\\
$80\%$ thể tích bể: $6\,750 \times 80\% = 5\,400\text{ l}$.\\
Thời gian chảy: $5\,400 : 45 = 120\text{ (phút)} = 2\text{ giờ}$.\\
\textbf{Đáp án của câu hỏi là D}
}
\end{ex}

\textbf{Phần 2. Trắc nghiệm đúng sai}
\setcounterpageref{ex}{0}

\begin{ex} %[5H8H48]
Xét tính đúng sai của các biểu thức so sánh thể tích:
\choiceTF
{\True $3\text{ m}^3\; 15\text{ dm}^3 = 3{,}015\text{ m}^3$.}
{$4{,}5\text{ m}^3 < 4\,500\text{ dm}^3$.}
{\True $2\text{ dm}^3\; 50\text{ cm}^3 > 2{,}005\text{ dm}^3$.}
{\True $1\,500\,000\text{ cm}^3 = 1{,}5\text{ m}^3$.}
\loigiai{
\begin{itemize}
    \item $3 + 15 : 1000 = 3{,}015\text{ m}^3$. Đúng.
    \item $4{,}5\text{ m}^3 = 4500\text{ dm}^3$. Dấu $<$ là Sai.
    \item $2{,}05\text{ dm}^3 > 2{,}005\text{ dm}^3$. Đúng.
    \item $1500000 : 1000000 = 1{,}5\text{ m}^3$. Đúng.
\end{itemize}
\begin{itemize}
    \item a) \textbf{Đúng}.
    \item b) \textbf{Sai}.
    \item c) \textbf{Đúng}.
    \item d) \textbf{Đúng}.
\end{itemize}
}
\end{ex}

\begin{ex} %[5H8V48]
Một khối gỗ hình hộp chữ nhật có chiều dài $1{,}5\text{ m}$, chiều rộng $0{,}8\text{ m}$ và chiều cao $0{,}5\text{ m}$. Biết mỗi đề-xi-mét khối gỗ nặng $0{,}7\text{ kg}$. Xét tính đúng sai:
\choiceTF
{\True Thể tích khối gỗ là $0{,}6\text{ m}^3$.}
{\True Thể tích khối gỗ bằng $600\text{ dm}^3$.}
{Khối lượng của khối gỗ là $42\text{ kg}$.}
{\True Khối lượng của khối gỗ là $420\text{ kg}$.}
\loigiai{
\begin{itemize}
    \item Thể tích: $1{,}5 \times 0{,}8 \times 0{,}5 = 0{,}6\text{ m}^3$. Đúng.
    \item $0{,}6\text{ m}^3 = 600\text{ dm}^3$. Đúng.
    \item Khối lượng: $600 \times 0{,}7 = 420\text{ kg} \neq 42\text{ kg}$. Ý 3 Sai, Ý 4 Đúng.
\end{itemize}
\begin{itemize}
    \item a) \textbf{Đúng}.
    \item b) \textbf{Đúng}.
    \item c) \textbf{Sai}.
    \item d) \textbf{Đúng}.
\end{itemize}
}
\end{ex}

\textbf{Phần 3. Trả lời ngắn}
\setcounterpageref{ex}{0}

\begin{ex} %[5H8H48]
Đổi $4\text{ m}^3\; 65\text{ dm}^3$ ra đề-xi-mét khối.
\shortans{4065}
\loigiai{
\textbf{Đáp số:} $4065$\\
$4 \times 1\,000 + 65 = 4\,065\text{ dm}^3$.
}
\end{ex}

\begin{ex} %[5H8H48]
Đổi $750\,000\text{ cm}^3$ ra mét khối (dùng dấu phẩy nếu có).
\shortans{0,75}
\loigiai{
\textbf{Đáp số:} $0{,}75$\\
$750\,000 : 1\,000\,000 = 0{,}75\text{ m}^3$.
}
\end{ex}

\begin{ex} %[5H8V48]
Một bể bơi dạng hình hộp chữ nhật có kích thước dài $15\text{ m}$, rộng $6\text{ m}$, sâu $1{,}5\text{ m}$. Tính thể tích nước tối đa mà bể có thể chứa theo mét khối.
\shortans{135}
\loigiai{
\textbf{Đáp số:} $135$\\
$15 \times 6 \times 1{,}5 = 135\text{ m}^3$.
}
\end{ex}

\begin{ex} %[5H8V48]
Một thùng chứa $1{,}8\text{ m}^3$ dầu. Người ta chia đều vào các can $20\text{ l}$. Hỏi chia được bao nhiêu can dầu như thế?
\shortans{90}
\loigiai{
\textbf{Đáp số:} $90$\\
$1{,}8\text{ m}^3 = 1\,800\text{ l} \Rightarrow 1\,800 : 20 = 90\text{ can}$.
}
\end{ex}

\begin{ex} %[5H8V48]
Một bể nước hình hộp chữ nhật có diện tích đáy $3\text{ m}^2$, chiều cao $1{,}2\text{ m}$. Bể đang chứa $2\,700\text{ l}$ nước. Hỏi cần đổ thêm bao nhiêu lít nước nữa để đầy bể?
\shortans{900}
\loigiai{
\textbf{Đáp số:} $900$\\
Thể tích bể: $3 \times 1{,}2 = 3{,}6\text{ m}^3 = 3\,600\text{ l}$.\\
Cần thêm: $3\,600 - 2\,700 = 900\text{ l}$.
}
\end{ex}

\begin{ex} %[5H8C48]
Một bể hình hộp chữ nhật có chiều dài $2\text{ m}$, chiều rộng $1{,}5\text{ m}$ và mực nước đang cao $0{,}8\text{ m}$. Người ta thả vào bể một hòn non bộ thì mực nước trong bể dâng lên cao $0{,}95\text{ m}$ (hòn non bộ chìm hoàn toàn). Tính thể tích hòn non bộ theo đề-xi-mét khối.
\shortans{450}
\loigiai{
\textbf{Đáp số:} $450$\\
Nước dâng thêm: $0{,}95 - 0{,}8 = 0{,}15\text{ m}$.\\
Thể tích hòn non bộ: $2 \times 1{,}5 \times 0{,}15 = 0{,}45\text{ m}^3 = 450\text{ dm}^3$.
}
\end{ex}
